% E68d, 26 November 2026 UTC. Drafted for Edward B. Baker III with
% OpenAI Codex assistance. Actual serving model and effort unavailable.
% Input descent generation d9a6557377462fd3, assessed in the E63--E67
% receipt. Exact-band and supplied marked round alternatives only.

\subsection{Chiral overlap on an exact comparison disk}
\label{sec:exact-overlap}

For the two top Weyl eigenvalues \(z_\pm\), define the conformally
invariant measure \(d\mathcal O=\min(z_+^2,z_-^2)dV\).
Fix a strict energy band \(e\le 2T-\delta\), an integer \(k\ge2\)
and \(\varepsilon>0\). There is \(a>0\), depending on these three
parameters, such that every volume-one Yamabe representative \(g\)
in the nonstandard exact band of Theorem~\ref{thm:exact-spectral-gap}
has an oriented embedding \(F\) of a neighborhood of the Euclidean
closed ball \(\overline B_2\) and a scale \(r>0\) with
\[
 \|r^{-2}F^*g-\delta_{\mathrm{Euc}}\|_{C^k(\overline B_2)}
       <\varepsilon,\qquad
 \mathcal O(F(B_1))\ge a .
\]
For each fixed sublevel \(\Delta_{\mathrm{sp}}\le m_\delta+b\)
with \(0\le b<6\kappa_\delta\), one can also require \(r\) to
have a uniform positive lower bound. The disk in this assertion is
selected by the exact comparison.

To prove the assertion, consider a countersequence along which the
suprema over such charts tend to zero. The exact compactness
alternative used in Theorem~\ref{thm:exact-spectral-gap} gives a
smooth closed limit or a marked smooth ALE limit with a positive
exact same-type completion. In each case the closed limit class
has nonzero energy in both chiralities by the closed gap and
signature. A smooth Bach-flat four-metric has zero obstruction
tensor. It admits local analytic conformal representatives
\cite[Theorem~1.3(c)]{LassasLiimatainen2019ConformalHarmonic}.
In these representatives each chiral Weyl block is analytic,
so vanishing on an open set propagates by the analytic identity
theorem through overlapping charts. Hence each nonzero locus is
open dense, and their intersection contains a small normal chart
with positive overlap capture. For the ALE branch take this chart
away from the filled point and use the metric before completion.
In either branch choose the normal scale with error below
\(\varepsilon/2\) on \(\overline B_2\). Its whole fixed
compact chart transfers to the late sources. For the ALE branch
the source scale includes the constant factor of the pointed
rescaling. Smooth convergence and continuity of the top
eigenvalues give a strictly positive limiting overlap on the
transferred disk, contrary to the countersequence. Halving the
positive infimum of the suprema gives one suitable \(a\).
In the strict deficit sublevel the ALE branch is excluded.
In a closed subsequential limit the full-band supremum is at
least \(2a\). Choose there a chart with capture strictly above
\(a\) and a fixed positive scale. Its transfer rules out a
sequence with no chart of scale \(r\ge i^{-1}\) and capture
at least \(a\), giving the joint lower bounds.

Writing \(D_\pm=\int(\lambda_1^\pm-\lambda_2^\pm)^2dV\) as in
Section~\ref{sec:endpoint-charge}, trace freeness gives
\[
 Y^2=48e-36\mathcal O(M)-12(D_++D_-)-\Delta_{\mathrm{sp}} .
\]
Thus in the exact band \(Y^2\le48e-\Delta_{\mathrm{sp}}-36a\).
In particular no sequence in that band has vanishing chiral
overlap. The bound holds at a class attaining \(m_\delta\).
Extremality there is through exact classes and supplies no
unrestricted variation of the spectral deficit.

\subsection{Volume and curvature at the round exit}
\label{sec:round-joint-exit}

Suppose \(g_i\) are volume-one positive Yamabe metrics on one
nonstandard oriented smooth homotopy four-sphere with
\(e_i\to I^+(M)\), \(e_i\le2T-\delta\), and
\(\Delta_{\mathrm{sp}}([g_i])\to0\). Suppose the whole
marked filled-base comparison in
Theorem~\ref{prop:local-min} and the subsequent cap argument
is given, including its strong first-order convergence and local
full \(L^2\) curvature bound. The spectral transfer makes the
filled base \((N,h)\) conformally round. Denote its finitely many
ordinary marks by \(P\), and by \(s\) the single exterior label
retaining both limiting chiral charges. Let \(q_i\) be the
unit-volume modified Yamabe metric in \([g_i]\).

\begin{proposition}
\label{prop:round-joint-exit}
In these comparisons both \(g_i\) and \(q_i\) have volume tending
to zero on every fixed compared compact. At each fixed small
cutting radius, the actual exterior at \(s\) carries all limiting
volume in both gauges, both chiral integrals and the spectral
integral. If the comparison metrics converge smoothly on
compacts, this same exterior also carries all limiting full
curvature in both gauges.
\end{proposition}

\begin{proof}
On a subsequence put \(Y_i\to y\) and \(S=8\sqrt6\pi\).
The sharp spectral bound \(\sigma_i\le4e_i/3\) and the strict
energy band give
\[
 0<48\delta\le y^2=36\lim_i\sigma_i\le48I^+(M)\le S^2-48\delta .
\]
On the comparison write \(\Phi_i^*g_i=u_i^2m_i\).
Use \(\mathcal L_m=-6\Delta_m+R_m\) and \(Q_m\) for its
quadratic form, with \(\Delta_m=\operatorname{div}_m\nabla_m\).
Strong first-order convergence bounds \(Q_{m_i}-Q_h\)
on each fixed compact by \(o(1)\) times the squared \(H^1\)
norm. Indeed the scalar density is the divergence of a term
linear in the first metric derivatives, plus a term quadratic
in them, with smooth metric coefficients. From the Yamabe
equation and a fixed squared cutoff one obtains
\[
 Q_{m_i}(\chi u_i)
   =Y_i\int \chi^2u_i^4dV_{m_i}
       +6\int u_i^2|d\chi|_{m_i}^2dV_{m_i} .
\]
Hence \(u_i\) is locally bounded in \(H^1\), and any weak
limit \(u\ge0\) solves \(\mathcal L_hu=yu^3\) with
\(\|u\|_{4,h}\le1\). The equation extends across \(P\).
Testing first with squared puncture cutoffs gives finite
energy there, since the \(L^4\) norms of their gradients are
bounded and the \(L^4\) tails of \(u\) vanish. The points
then have zero \(H^1\) capacity. Testing by \(u\) gives
\(S\|u\|_4^2\le y\|u\|_4^4\), which forces \(u=0\).
In the cutoff identity, the nonlinear term is bounded by
\((y+o(1))\|\chi u_i\|_{4,m_i}^2\) by unit volume and
Hölder. On small fixed coordinate balls the sharp Sobolev
inequality has coefficient \(S-o(1)>y\), with the error made
small by the chart radius
\cite[Theorems 3.2 and 3.3]{LeeParker1987Yamabe}.
The \(L^2\) cutoff terms tend to zero by local Rellich convergence.
Absorption gives strong local \(H^1\cap L^4\) convergence
to zero.

On the actual exits, cutoffs which transition on fixed
compared annuli now have physical squared gradient integral
tending to zero by Hölder and conformal invariance of the
fourth-power gradient integral. The source Yamabe inequality
forces \(y\sqrt{\theta_j}\le y\theta_j\) for each limiting
exit fraction \(\theta_j\). Their sum is one, so one exit has
unit limiting volume. Fix the cap radius and buffer, and place
the cutoff transition deeper into that raw exit than the
innermost splice. Its entire support compares conformally on the
unchanged part of the cap, giving a Yamabe upper limit at most
\(y\). If this exit were a label other than \(s\), its positive
cap would have vanishing chiral energy and hence Yamabe constant
tending to \(S\) by Gauss--Bonnet. Thus it is \(s\).
Take the source limit at these fixed choices before the cap
exhaustion. Changing between fixed small cutting radii
only adds or removes an evacuated compared annulus.

For \(q_i\), put \(\beta_i=R_{q_i}-6L_{q_i}\le0\), the modified
Yamabe value
\cite[Lemma~2.1 and Theorem~1.7]{ChenZhu2014Conformal}.
The chain \(Y_i\le\int R_{q_i}dV_{q_i}
   \le\beta_i+6\sqrt{\sigma_i}\) implies
\(\beta_i\to0\) and \(R_{q_i}-Y_i\to0\) in \(L^2(q_i)\).
It also gives \(L_{q_i}\to y/6\) in \(L^2(q_i)\), by
equality in the limiting unit-volume Cauchy--Schwarz inequality.
Consequently \(L^2dV-(y^2/36)dV_{q_i}\to0\) in total variation.
Both chiral masses vanish locally in the round comparison and
on the other exteriors. They bound the invariant spectral
measure, so the modified volume is also evacuated from these
sets and concentrates at \(s\).

Assume now \(m_i\to h\) smoothly on compacts. Write
\(\Phi_i^*q_i=v_i^2m_i\).
The ordinary factor \(u_i\) converges locally uniformly to zero
by local Sobolev iteration in its equation, using the strong
\(L^4\) convergence to absorb the \(L^2\) potential \(Y_i u_i^2\).
The modified equation is
\[
 \mathcal L_{m_i}v_i=(6L_{m_i}+\beta_i v_i^2)v_i .
\]
Since \(\beta_i\le0\), \(v_i\) is a positive subsolution of
\(-6\Delta_{m_i}v_i\le C v_i\) on buffered compacts, with a
uniform local \(C\). Its \(L^4\) convergence to zero then gives
local uniform convergence. At a fixed \(x_0\in N-P\), normalize
the two factors by their values at \(x_0\). Harnack and interior
elliptic estimates give subsequential \(C^2_{\mathrm{loc}}\)
limits, positive and \(\mathcal L_h\)-harmonic
\cite[Sections 6.1 and 8.5--8.8]{GilbargTrudinger2001Elliptic}.
For this regularity the normalized factors first have local
\(W^{2,p}\) bounds for finite \(p\), since their linear potentials
are bounded. The top eigenvalues in smoothly moving frames
are locally Lipschitz with bounded constants. Together with the
\(W^{2,p}\) bounds, the equations then give interior
\(C^{2,\alpha}\) bounds for \(0<\alpha<1\).

To identify the positive limits, take a cutoff equal to one
away from \(s\) and turning off on a fixed compared annulus in
that exit. Its support has vanishing volume in both gauges.
In the modified gauge, the source Yamabe inequality implies
\[
 \left(Y_i-\|R_{q_i}-Y_i\|_{2,q_i}
  -Y_i\sqrt{\operatorname{Vol}_{q_i}(\operatorname{supp}\chi)}\right)
 \|\chi\|_{4,q_i}^2
 \le6\int|d\chi|_{q_i}^2dV_{q_i}=O(v_i(x_0)^2).
\]
The ordinary test replaces \(v_i(x_0)\) by \(u_i(x_0)\)
and has zero scalar error. The positive lower coefficient
bounds the volume after the corresponding constant rescaling
on all the parts where the cutoff is one. After interior
convergence and exhaustion, either limiting factor is locally
\(L^4\) across the other marks. In flat conformal coordinates
it extends there as a harmonic function after multiplication
by the conformal factor of \(h\). To see removability, test
against a smooth function times a shrinking cutoff. The
second-derivative error has \(L^{4/3}\) norm \(O(r)\), so the
weak Laplace equation passes through the mark. Stereographic
projection from \(s\) now gives a positive harmonic function
on \(\mathbb R^4\). It is constant, with its value fixed by
normalization. Both limits therefore give the same flat
rescaled metric on the compared punctured base. The squared
curvature integrals on its fixed compared compacts tend to zero
by convergence and constant-scale invariance.

Finally, at a fixed small face around any \(j\ne s\), the
entire actual exterior has \(\chi=1\), since attaching its
standard ball cap gives a homotopy four-sphere. Its squared
scalar and Weyl integrals vanish by the volume and chiral
accounting above. The Euler boundary term converges to
\(8\pi^2\): the comparison after constant rescaling is \(C^2\)
and flat across the filled mark, and the outward normal from
that exterior points away from the puncture, as for a filled
disk. The boundary term is invariant under constant scaling
and continuous under \(C^2\) convergence on the collar.
In the convention \(8\pi^2\chi\) this is the boundary
integral of the flat filled disk
\cite[Section 7, (7.4)]{TianViaclovsky2005ALE}.
Gauss--Bonnet on the entire exterior now gives
\(\int|Z|^2dV\to0\). The same formula for the
\(C^{2,\alpha}\) modified metric follows by \(C^2\) smooth
approximation on each fixed source. The full curvature
decomposition gives the asserted allocation in both gauges.
\end{proof}

The common exit and its flat round-side model control the limiting
accounting. A strict restart comparison still requires the
finite-source difference in \eqref{eq:signed-restart} to exceed
the minimizing excess. The exact graft of
Proposition~\ref{prop:graft} has equality in that difference.
